\documentclass[11pt]{amsart}
\usepackage{amsmath,amssymb,amsthm}
\usepackage[margin=1.05in]{geometry}

\numberwithin{equation}{section}
\newtheorem{theorem}{Theorem}[section]
\newtheorem{proposition}[theorem]{Proposition}
\newtheorem{lemma}[theorem]{Lemma}
\newtheorem{corollary}[theorem]{Corollary}
\theoremstyle{definition}
\newtheorem{definition}[theorem]{Definition}
\theoremstyle{remark}
\newtheorem{remark}[theorem]{Remark}
\newtheorem{example}[theorem]{Example}

\newcommand{\Z}{\mathbb{Z}}
\newcommand{\Q}{\mathbb{Q}}
\newcommand{\R}{\mathbb{R}}
\newcommand{\N}{\mathbb{N}}
\newcommand{\C}{\mathbb{C}}
\newcommand{\F}{\mathbb{F}}
\newcommand{\HH}{\mathbb{H}}
\newcommand{\cH}{\mathcal{H}}
\newcommand{\cT}{\mathcal{T}}
\newcommand{\cM}{\mathcal{M}}
\newcommand{\Nrd}{\operatorname{Nrd}}
\newcommand{\Ns}{\mathrm{N}}
\newcommand{\pp}{\mathfrak{p}}
\newcommand{\id}{\mathrm{id}}

\PassOptionsToPackage{hyphens}{url}
\usepackage[hidelinks,bookmarks=false]{hyperref}

\title[The Ramanujan route to (C4) is closed]{Condition \textup{(C4)} and the nested/partition dichotomy:\\ why Ramanujan bounds cannot decide\\ non-amenability of the Hurwitz groupoid}
\author{Ruqing Chen}
\address{GUT Geoservice Inc., Montreal, Canada}
\email{ruqing@hotmail.com}
\thanks{Sources, code and data for the entire series: \url{https://github.com/Ruqing1963/localized-bost-connes}; this paper is in \texttt{papers/XI-C4/}.}
\date{}

\begin{document}

\begin{abstract}
Condition (C4) of Paper X of the series asks whether the measured groupoid of the Hurwitz
Bost--Connes system is non-amenable, equivalently whether the associated factor escapes the
injectivity that collapsed every earlier model (Paper IX). We examine the natural route ---
Lubotzky--Phillips--Sarnak generators \cite{LPS}, the Bruhat--Tits tree, Deligne's bound
\cite{Deligne}, a spectral gap for a Hecke--Markov operator in the sense of Kesten
\cite{Kesten} --- and show that it is closed, for reasons that are structural rather than
technical.

The obstruction is a dichotomy. A $(p+1)$-regular tree adjacency operator requires the $p+1$
branches at a vertex to \emph{partition}; but the Ore and LCM conditions of Paper X force
the range projections to be \emph{nested}, $e_je_k=e_{a_j\vee a_k}=e_p\ne0$ --- the central
element $p$ lies in every maximal right ideal of norm $p$ --- with $\varphi(e_je_k)=p^{-2\beta}>0$
and $\sum_k\varphi(e_{a_k})=(p+1)p^{-\beta}\ne1$. The very property that keeps the
high-temperature region alive --- absence of a Cuntz relation --- is the property that
destroys the tree.

Forcing the partition does not help, and the reason is sharper than a numerical mismatch: in
any unital algebra generated by isometries $v_a$ with $v_av_b=v_{ab}$, orthogonality of two
range projections of norm $p$ forces $e_p=0$ and hence $1=v_p^*v_p=0$. \emph{The boundary
quotient of the Hurwitz system is the zero algebra.} Even a partition in mean,
$\sum_k\varphi(e_{a_k})=1$, would force $\beta=\log(p+1)/\log p$, a value that varies with $p$
($1.2619,\,1.1133,\,1.0686,\dots$ for $p=3,5,7$). Tree structure and thermodynamics are
mutually exclusive here.

Two further points are recorded. The proposed operator
$\cM_p=\frac{1}{p+1}\sum_kv_{a_k}^*(\cdot)v_{a_k}$ is unital completely positive but does not
preserve $\varphi$: one has $\varphi\circ\cM_p\le p^{\beta}\varphi$ with the naive bound
$\|\cM_p\|_{L^2(\varphi)}\le p^{\beta/2}>1$, so no contraction is available, let alone a gap
below $2\sqrt p/(p+1)$. And the statement of the spectral gap is not well posed: for a
non-singular action the Koopman representation satisfies $\pi(a)1=N(a)^{\beta/2}1_{a\Omega}\ne1$,
so $\C1$ is not invariant and $L^2_0=L^2\ominus\C1$ is not a subrepresentation.

We close with the reformulation we believe is correct: by Takesaki duality the factor is
injective if and only if its continuous core is, so (C4) is equivalent to non-amenability of
the \emph{Maharam extension}, which is measure preserving. That does not settle it --- non-amenable groups do have
amenable infinite-measure-preserving actions --- but it places the question in a
measure-preserving setting, where the obstruction identified here does not apply.
Condition (C4) has since been settled, in the negative, in Paper XXVIII of the series,
by the properness route rather than any spectral one, confirming the warning of this paper.
\end{abstract}

\maketitle

\section{What the route requires}

We keep the notation of Paper X: $\cH$ the Hurwitz order, $P=\cH\setminus\{0\}$,
$G=\HH(\Q)^\times$, $\Nrd$ the reduced norm, $\cT(P)$ the Nica--Toeplitz algebra with
isometries $v_a$, projections $e_a=v_av_a^*$, dynamics $\sigma_t(v_a)=\Nrd(a)^{it}v_a$, and
$\varphi$ an extremal $\mathrm{KMS}_\beta$ state with $\varphi(e_a)=\Nrd(a)^{-\beta}$.

For an odd prime $p$ let $a_1,\dots,a_{p+1}$ represent the $p+1$ principal right ideals of
reduced norm $p$, which are the $p+1$ neighbours of a vertex in the $(p+1)$-regular
Bruhat--Tits tree of $SL_2(\Q_p)$, $\HH\otimes\Q_p$ being split for $p$ odd (Paper X).

\emph{The series.} The foundational paper is \cite{Main}; Papers II, III, V, IX and X, cited
by numeral, are in preparation, and the facts we use from them are re-proved where they
are used.

The route under examination is: (i) form
$\cM_p=\frac{1}{p+1}\sum_kv_{a_k}^*(\cdot)v_{a_k}$; (ii) identify it with a tree adjacency
operator; (iii) invoke Deligne's bound to get
$\|\cM_p\|_{L^2_0}\le2\sqrt p/(p+1)<1$; (iv) conclude non-amenability by a Kesten-type
criterion. We show that (ii) and (iii) both fail, and that repairing (ii) destroys the
system.

\section{The projections do not partition}

\begin{proposition}[Nested, not orthogonal]\label{prop:nested}
Let $j\ne k$. Then
\[
e_{a_j}e_{a_k}=e_{a_j\vee a_k}\ne0,\qquad
\varphi\bigl(e_{a_j}e_{a_k}\bigr)=p^{-2\beta}>0,
\qquad
\sum_{k=1}^{p+1}\varphi\bigl(e_{a_k}\bigr)=(p+1)\,p^{-\beta}.
\]
In particular $\{e_{a_k}\}$ is not an orthogonal family and $\sum_ke_{a_k}$ is not a scalar,
so it is not a partition of unity for any $\beta$.
\end{proposition}

\begin{proof}
In the Nica--Toeplitz algebra $e_ae_b$ is the projection onto $aP\cap bP$, which is
$(a\vee b)P$ by the LCM property, nonempty by the Ore property (Paper X). For distinct right
ideals of reduced norm $p$ one has $a_j\cH\cap a_k\cH=p\cH$: the central element
$p=a_j\bar a_j=a_k\bar a_k$ lies in both, the two ideals are distinct maximal right ideals
containing $p\cH$, so $a_j\cH+a_k\cH=\cH$ and $[\cH:a_j\cH\cap a_k\cH]=[\cH:a_j\cH][\cH:a_k\cH]=p^4=[\cH:p\cH]$.
Hence $a_j\vee a_k=p$ up to units, $e_{a_j}e_{a_k}=e_p$ and $\varphi(e_p)=\Nrd(p)^{-\beta}=p^{-2\beta}$.
The last identity is $\varphi(e_{a_k})=p^{-\beta}$ summed over $p+1$ terms.
\end{proof}

\begin{corollary}[The tree picture is unavailable]\label{cor:notree}
The adjacency operator of a $(p+1)$-regular tree acts on $\ell^2$ of the vertex set by summing
over $p+1$ \emph{disjoint} branches, and its analysis --- Kesten's $2\sqrt p$, Deligne's
bound for the LPS quotients --- rests on that disjointness. By
Proposition~\ref{prop:nested} the branches of $\cT(P)$ overlap, with overlap mass
$p^{-2\beta}$ at every pair. No identification of $\cM_p$ with a tree adjacency operator is
therefore available, and the Ramanujan bound has no object to bound.
\end{corollary}

\begin{remark}[Numerical illustration]\label{rem:num1}
At $p=3$: $\sum_k\varphi(e_{a_k})=1.3333$ and $\varphi(e_je_k)=0.1111$ for $\beta=1$;
$0.4444$ and $0.01235$ for $\beta=2$. At $p=13$, $\beta=2$: $0.0828$ and $3.5\cdot10^{-5}$.
The overlaps shrink as $\beta$ grows but never vanish, and the branch masses never sum to $1$.
\end{remark}

\begin{remark}[This is exactly the property we wanted]\label{rem:tension}
The failure is not incidental. The Ore and LCM conditions were imposed in Paper X precisely
to guarantee that no Cuntz relation holds, which is what keeps the high-temperature region
$\beta<2$ from being excluded, as it is in the affine systems of Paper II and the free
variant of Paper V. Nestedness is the feature that makes the Hurwitz model viable; and
nestedness is the feature that forbids the tree. The two requirements are in direct
opposition.
\end{remark}

\section{The Markov operator does not contract}

\begin{proposition}[$\cM_p$ is unital but not state-preserving]\label{prop:notpreserving}
$\cM_p(x)=\frac{1}{p+1}\sum_kv_{a_k}^*xv_{a_k}$ is unital and completely positive, and
\begin{equation}\label{eq:notpres}
\varphi\bigl(\cM_p(x)\bigr)=\frac{p^{\beta}}{p+1}\,\varphi\Bigl(x\sum_{k}e_{a_k}\Bigr)
\qquad(x\in\cT(P)).
\end{equation}
Since $\sum_ke_{a_k}$ is not a scalar by Proposition~\ref{prop:nested}, $\varphi\circ\cM_p\ne\varphi$.
Consequently $\cM_p$ induces on $L^2(\cT(P),\varphi)$ only the bound
\[
\bigl\|\cM_p\bigr\|_{L^2(\varphi)}\ \le\ p^{\beta/2},
\]
which exceeds $1$ for every $\beta>0$, and no contraction --- still less a gap below
$2\sqrt p/(p+1)$ --- follows from these relations.
\end{proposition}

\begin{proof}
Unitality is $\frac1{p+1}\sum_kv_{a_k}^*v_{a_k}=1$; complete positivity is clear. The KMS
condition with the entire element $v_a^*$, $\sigma_{i\beta}(v_a^*)=\Nrd(a)^{\beta}v_a^*$,
gives $\varphi(v_a^*xv_a)=\Nrd(a)^{\beta}\varphi(xe_a)$, the scaling identity of Paper III,
whence \eqref{eq:notpres} with $\Nrd(a_k)=p$. For the bound: $e_a$ is fixed by $\sigma$, so
$\varphi(e_ay)=\varphi(ye_a)$ for all $y$, whence $\varphi(xe_a)=\varphi(e_axe_a)$ and, for
$x\ge0$, $\varphi(x)=\varphi(e_axe_a)+\varphi(e_a^\perp xe_a^\perp)\ge\varphi(e_axe_a)$; so
$\varphi\circ\cM_p\le p^{\beta}\varphi$ on positive elements, and a unital completely positive
map $\Phi$ with $\varphi\circ\Phi\le C\varphi$ satisfies
$\varphi(\Phi(x)^*\Phi(x))\le\varphi(\Phi(x^*x))\le C\varphi(x^*x)$ by the Kadison--Schwarz
inequality, i.e.\ extends to $L^2(\varphi)$ with norm at most $C^{1/2}$.
\end{proof}

\begin{remark}\label{rem:sameas}
Proposition~\ref{prop:notpreserving} is the Hurwitz instance of the phenomenon recorded in
Paper III (its Proposition 5.1): the transfer operator $x\mapsto v_a^*xv_a$ never preserves a
$\mathrm{KMS}_\beta$ state, because the KMS condition converts it into
$\Nrd(a)^\beta\varphi(\cdot\,e_a)$. Whatever spectral theory one wishes to do with $\cM_p$
must first address this, and normalizing by $(p+1)^{-1}$ does not.
\end{remark}

\section{The spectral gap statement is not well posed}

\begin{proposition}[$\C1$ is not invariant]\label{prop:notinvariant}
For the non-singular action of $G$ on $(\Omega,\nu_\beta)$ with
$\nu_\beta(a\,\cdot)=\Nrd(a)^{-\beta}\nu_\beta$, the Koopman representation is
\[
\bigl(\pi(a)\xi\bigr)(x)=\Nrd(a)^{\beta/2}\,\xi\bigl(a^{-1}x\bigr),
\qquad\text{so}\qquad
\pi(a)1=\Nrd(a)^{\beta/2}\,\mathbf 1_{a\Omega}\ \ne\ 1
\]
for $a\notin P^\times$. Hence $\C1$ is not a $G$-invariant subspace, $L^2_0=L^2\ominus\C1$ is
not a subrepresentation, and the quantity $\|\cM_p|_{L^2_0}\|$ is not the norm of an operator
on an invariant subspace.
\end{proposition}

\begin{proof}
Immediate from the Radon--Nikodym cocycle; $\mathbf 1_{a\Omega}\ne1$ because $a\Omega$ is a
proper clopen subset, of measure $\Nrd(a)^{-\beta}<1$.
\end{proof}

\begin{remark}[Why this matters]\label{rem:whymatters}
For a \emph{measure-preserving} action the decomposition $L^2=\C1\oplus L^2_0$ is invariant
and a spectral gap on $L^2_0$ is meaningful and does obstruct amenability. That is the
setting of Kesten's criterion and of the standard argument ``non-amenable group acting
measure-preservingly $\Rightarrow$ non-amenable action''. Neither applies here: our measure is
quasi-invariant with a nontrivial cocycle, which is the whole point of a
$\mathrm{KMS}_\beta$ state, and it is precisely for non-singular actions that non-amenable
groups routinely admit amenable actions \cite{AD,Zimmer} --- the reason (C4) was posed as
open in Paper X rather than settled there.
\end{remark}

\section{Forcing the partition destroys the system}

\begin{theorem}[Forcing the partition kills the algebra]\label{thm:noboundary}
Let $A$ be a unital $C^*$-algebra generated by isometries $v_a$, $a\in P$, with
$v_av_b=v_{ab}$, and let $p$ be an odd prime. If two range projections $e_{a_j}$, $e_{a_k}$,
$j\ne k$, of reduced norm $p$ are orthogonal in $A$, then $A=0$. In particular the
``boundary quotient'' $\partial\cT(P)$ of $\cT(P)$ by the relations
$\sum_{k=1}^{p+1}e_{a_k}=1$, for any single odd prime $p$, is the zero algebra, and carries
no state, $\mathrm{KMS}$ or otherwise.
\end{theorem}

\begin{proof}
Since $\Nrd(a_j)=p$ we have $p=a_j\bar a_j$ with $\bar a_j\in P$, so
$v_p=v_{a_j}v_{\bar a_j}$ and $e_p=v_pv_p^*=v_{a_j}e_{\bar a_j}v_{a_j}^*\le v_{a_j}v_{a_j}^*=e_{a_j}$;
likewise $e_p\le e_{a_k}$. If $e_{a_j}e_{a_k}=0$ then $e_p\le e_{a_j}\wedge e_{a_k}=0$, so
$v_p=e_pv_p=0$ and $1=v_p^*v_p=0$. A sum of projections equal to $1$ is an orthogonal sum:
$e_j=e_j\sum_ke_k$ gives $\sum_{k\ne j}e_je_ke_j=0$ with positive summands, so
$e_je_ke_j=0$ and $e_je_k=0$. Hence the relation $\sum_ke_{a_k}=1$ already implies $1=0$.
\end{proof}

\begin{remark}[Partition in mean]\label{rem:betavals}
One might weaken the partition to a partition \emph{in mean}, asking only that
$\sum_k\varphi(e_{a_k})=1$ for a $\mathrm{KMS}_\beta$ state of $\cT(P)$ itself. This forces
$(p+1)p^{-\beta}=1$, i.e.
\[
\beta=\beta(p):=\frac{\log(p+1)}{\log p},
\]
and $p\mapsto\beta(p)$ is strictly decreasing on $(1,\infty)$, since
$\frac{d}{dp}\log\bigl(\log(p+1)/\log p\bigr)$ is negative there: $\beta(3)=1.261860$,
$\beta(5)=1.113283$, $\beta(7)=1.068622$, $\beta(11)=1.036287$, $\beta(13)=1.028893$,
$\beta(101)=1.002135$, $\beta(1009)=1.000143$, decreasing to $1$. Any two primes already
give a contradiction, so not even the mean version of the partition is compatible with a
$\mathrm{KMS}_\beta$ state at more than one prime.
\end{remark}

\begin{corollary}[The dichotomy]\label{cor:dichotomy}
For the Hurwitz system the following are mutually exclusive:
\begin{enumerate}
\item the range projections at some prime partition the unit, so that tree and
Ramanujan methods apply;
\item the algebra is nonzero, in particular a $\mathrm{KMS}_\beta$ state exists.
\end{enumerate}
Consequently no argument that first passes to a partition can decide (C4), the object it
would decide being zero.
\end{corollary}

\begin{remark}[Contrast with the classical $ax+b$ case]\label{rem:contrast}
In the affine systems of Paper II the Cuntz relation at $\pp$ reads
$\Ns\pp\cdot\Ns\pp^{-\beta}=1$, giving $\beta=1$ \emph{independently of $\pp$}, so a KMS state
survives at the single temperature $\beta=1$; there the $\Ns\pp$ translates of a fundamental
domain are genuinely disjoint, the acting semigroup being $ax+b$ rather than
multiplicative. Here the multiplicity is $p+1$ while the norm
is $p$, so the forced temperature $\log(p+1)/\log p$ depends on $p$ and no temperature
survives. The mismatch between multiplicity and norm --- $p+1$ against $p$, the discrepancy
being exactly the ramification-free part of the Bruhat--Tits degree --- is what makes the
Hurwitz boundary empty where the affine boundary is not.
\end{remark}

\section{The correct reformulation}

\begin{proposition}[Maharam]\label{prop:maharam}
(C4) is equivalent to non-amenability of the Maharam extension \cite{Maharam}
\[
G\curvearrowright\bigl(\widetilde\Omega\times\R,\ \tilde\nu_\beta\otimes e^{s}ds\bigr),
\qquad a\cdot(x,s)=\bigl(ax,\ s+\beta\log\Nrd(a)\bigr),
\]
on the dilation $\widetilde\Omega$ carrying the $G$-action and the extended measure
$\tilde\nu_\beta$, which \emph{preserves} the displayed infinite measure.
\end{proposition}

\begin{proof}
Invariance of the displayed measure: the pushforward under $a$ scales $\tilde\nu_\beta$ by
$\Nrd(a)^{-\beta}$ and $e^sds$ by $e^{\beta\log\Nrd(a)}$, which is the construction of the
Radon--Nikodym cocycle. For
the equivalence: amenability of a measured groupoid is equivalent to injectivity of its
von Neumann algebra \cite{AD}, and the von Neumann algebra of the Maharam extension is the
continuous core $M\rtimes_{\sigma^\varphi}\R$ of $M=\pi_\varphi(\cT(P))''$ \cite[Ch.~XII]{Takesaki}.
If $M$ is injective so is $M\rtimes_{\sigma}\R$, a crossed product by an amenable group; if
$M\rtimes_\sigma\R$ is injective then so is $M\otimes B(L^2(\R))\cong(M\rtimes_\sigma\R)\rtimes_{\hat\sigma}\R$
by Takesaki duality, hence $M$. So $M$ is injective iff the Maharam extension is amenable.
\end{proof}

\begin{remark}[Why this is progress, and why it is not a solution]\label{rem:maharamstatus}
The reformulation removes the obstruction of \S4: on the Maharam extension the measure is
invariant, so decompositions of $L^2$ into invariant pieces and spectral gap arguments are
meaningful. It does not resolve (C4). Amenability of an infinite-measure-preserving action of
a non-amenable group is possible --- the translation action of $G$ on itself with counting
measure is amenable as a groupoid, being principal and proper --- so non-amenability of
$G=\HH(\Q)^\times$ still does not transfer. What would be needed is a spectral gap for the
Koopman representation of the Maharam extension on the orthogonal complement of the invariant
functions, and we have not computed it.
\end{remark}

\section{Assessment}

\[
\begin{array}{lll}
\text{Step of the proposed route} & \text{Status} & \text{Reference}\\\hline
p+1\ \text{ideals}\ \leftrightarrow\ \text{tree neighbours} & \text{true} & \text{Paper X}\\
\HH(\Q)^\times\ \text{non-amenable} & \text{true, necessary} & \text{Paper X}\\
\text{the}\ e_{a_k}\ \text{partition} & \textbf{false: nested} & \text{Prop.~2.1}\\
\cM_p\ \text{a tree adjacency operator} & \textbf{false} & \text{Cor.~2.2}\\
\cM_p\ \text{preserves}\ \varphi & \textbf{false};\ \|\cM_p\|\le p^{\beta/2} & \text{Prop.~3.1}\\
L^2_0\ \text{a subrepresentation} & \textbf{false} & \text{Prop.~4.1}\\
\|\cM_p|_{L^2_0}\|\le2\sqrt p/(p+1) & \text{not well posed} & \text{Prop.~4.1}\\
\text{pass to the boundary to force a partition} & \textbf{the quotient is zero} & \text{Thm.~5.1}\\
\text{(C4) resolved} & \textbf{no} & \text{---}\\
\text{(C4)}\iff\text{Maharam non-amenable} & \text{true reformulation} & \text{Prop.~6.1}
\end{array}
\]

The route is closed, and Corollary~\ref{cor:dichotomy} says something stronger than that this
particular attempt fails: \emph{no} method requiring a partition of unity into $p+1$ pieces
can decide (C4), because the object on which such a partition exists is zero. This rules
out tree adjacency operators, random walks on the Bruhat--Tits building, and Ramanujan or
Kesten bounds in the form used above. The feature that makes them applicable is exactly the
Cuntz relation that the Ore and LCM conditions of Paper X were designed to exclude.

We are left with Proposition~\ref{prop:maharam} as the reformulation to pursue and with the
following as the concrete open questions. Does the Koopman representation of the Maharam
extension on the complement of the invariant functions have a spectral gap? If it does not,
(C4) fails and by Connes' classification, as in Paper IX, the Hurwitz factor collapses to
$R_\infty$, $R_\lambda$ or a Krieger factor like all the others, which would be a definitive
negative answer to the whole programme. If it does, one obtains the first non-injective
arithmetic Bost--Connes factor, and the deformation/rigidity machinery that Paper IX found
inapplicable for want of a spectral gap becomes actionable. Second: is there a variant of the Hurwitz
system in which the multiplicity $p+1$ and the norm $p$ are brought into agreement, so that a
boundary quotient with a nonempty KMS simplex exists? Remark~\ref{rem:contrast} suggests
looking for an order whose ideals of norm $\Ns\pp$ number exactly $\Ns\pp$; whether such an
order exists we do not know, and it is the one place where the dichotomy of
Corollary~\ref{cor:dichotomy} might be broken rather than circumvented.

\begin{thebibliography}{9}
\bibitem{AD} C. Anantharaman-Delaroche, J. Renault, \emph{Amenable Groupoids}, Monogr. Enseign. Math. 36, Geneva, 2000.
\bibitem{Deligne} P. Deligne, \emph{La conjecture de Weil. I}, Publ. Math. IHES 43 (1974), 273--307.
\bibitem{Kesten} H. Kesten, \emph{Symmetric random walks on groups}, Trans. Amer. Math. Soc. 92 (1959), 336--354.
\bibitem{LPS} A. Lubotzky, R. Phillips, P. Sarnak, \emph{Ramanujan graphs}, Combinatorica 8 (1988), 261--277.
\bibitem{Maharam} D. Maharam, \emph{Incompressible transformations}, Fund. Math. 56 (1964), 35--50.
\bibitem{Takesaki} M. Takesaki, \emph{Theory of Operator Algebras II}, Encyclopaedia Math. Sci. 125, Springer, 2003.
\bibitem{Zimmer} R. J. Zimmer, \emph{Ergodic Theory and Semisimple Groups}, Monogr. Math. 81, Birkh\"auser, 1984.
\bibitem{Main} R. Chen, \emph{KMS state uniqueness and phase transitions in localized Bost--Connes systems: from Chebotarev density to the Hardy--Littlewood conjecture}, preprint, 2026, DOI \href{https://doi.org/10.5281/zenodo.22152101}{10.5281/zenodo.22152101}.
\end{thebibliography}

\end{document}
